% Reissner-Mindlin curvature derivations -- typeset compilation of hand-derived notes.
% Compile with:  pdflatex RM_curvatures.tex   (or: tectonic RM_curvatures.tex)
% Requires the rendered sheets in ./img/sheet_00.png ... sheet_14.png
\documentclass[11pt]{article}
\usepackage[a4paper,margin=20mm]{geometry}
\usepackage{amsmath,amssymb}
\usepackage{graphicx}
\usepackage{xcolor}
\usepackage[hidelinks]{hyperref}
\usepackage{fancyhdr}
\definecolor{navy}{rgb}{0.12,0.20,0.45}
\pagestyle{fancy}\fancyhf{}\rhead{\small RM curvature derivations}\rfoot{\small\thepage}
\setlength{\headheight}{14pt}

\title{\color{navy}\bfseries Reissner--Mindlin Curvature Derivations\\
       \large\mdseries Typeset compilation of hand-derived notes}
\author{}
\date{Source: \texttt{RM derivations\_curvatures\_Handderived.pdf} (15 sheets)}

\begin{document}
\maketitle
\thispagestyle{fancy}

\section*{Scope}
MSG thin-walled \emph{Reissner--Mindlin} shell $\rightarrow$ beam model: closed-form
derivation of the shell curvature strains
$\kappa_{11},\ \kappa_{22},\ \kappa_{12}+\kappa_{21}$ and the transverse shears
$2\gamma_{13},\ 2\gamma_{23}$ in terms of the beam strains and the
rotation/displacement fluctuating functions, at a general reference surface.

\section*{Notation}
\begin{tabular}{@{}ll@{}}
$x_i,\ x_{i;\alpha}$ & position and its contour derivatives ($\dot{(\ )}=\partial/\partial\zeta_2$)\\
$k_{\alpha 1},k_{\alpha 2},k_{\alpha 3}$ & initial curvature components of the reference surface\\
$C^{ab}_{ij}=a_i\!\cdot\! b_j$ & direction-cosine matrix (shell $\leftrightarrow$ beam frames)\\
$\theta,\ \phi_3$ & beam rotation column and drilling rotation\\
$\omega_\beta$ & rotation fluctuating functions (RM extra DOFs)\\
$w_1,w_2,w_3$ & displacement fluctuating functions\\
$R_n=x_2\,x_{3;2}-x_3\,x_{2;2}$ & warping / enclosed-area kernel\\
$\kappa_1$ (twist),\ $\kappa_2,\kappa_3$ (bending) & beam generalized strains\\
\end{tabular}

\section*{Structure of the derivation}
The shell rotation vector is mapped to the beam frame through $C^{ab}$ and its
contour derivative $C^{ab}_{;\alpha}=-\tilde k^{s}_{\alpha}C^{ab}$. Eliminating the
drilling rotation $\phi_3$ (known from $\epsilon_{12}=\epsilon_{21}$) leaves the two
independent rotation fluctuations $\omega_\alpha$. The curvature strains follow as
first derivatives of the (rotation/displacement) fluctuating functions; e.g. the
transformation operator applied to the rotation gradient has the form
\[
\begin{bmatrix} 1 & 0 & 0\\[1pt] 0 & 1 & 0\\[1pt]
-\dfrac{C^{ab}_{31}}{C^{ab}_{33}} & -\dfrac{C^{ab}_{32}}{C^{ab}_{33}} & 1 \end{bmatrix}
\bigl(\theta'\,x_{1\kappa}+\varepsilon_\beta\,\omega_\beta'\,x_{1\kappa}
      +\varepsilon_\beta\,\omega_{\beta|1\kappa}\bigr),
\]
which, expanded, gives the $\kappa_{11},\kappa_{12}$ rows used throughout the sheets.

\medskip\noindent
\textbf{Faithful record.} Because the intermediate steps are index-dense, the exact
derivation is reproduced below as the original hand-derived sheets, embedded at high
resolution. (The equations above are transcribed from sheet~1; verify index detail
against the corresponding sheet.)

\newpage
% ---- embedded hand-derived sheets ----
\newcommand{\sheet}[2]{%
  \begin{center}{\color{navy}\bfseries Hand-derived sheet #1 / 15}\end{center}%
  \begin{center}\includegraphics[width=\linewidth,height=0.9\textheight,keepaspectratio]{img/sheet_#2.png}\end{center}%
  \clearpage}
\sheet{1}{00}\sheet{2}{01}\sheet{3}{02}\sheet{4}{03}\sheet{5}{04}
\sheet{6}{05}\sheet{7}{06}\sheet{8}{07}\sheet{9}{08}\sheet{10}{09}
\sheet{11}{10}\sheet{12}{11}\sheet{13}{12}\sheet{14}{13}\sheet{15}{14}
\end{document}
